\section{Discussion}

The empirical findings show that machine learning and deep learning models can provide useful support for Shariah-compliant stock price prediction and portfolio optimisation. The strongest forecasting performance is achieved by XGBoost, which records the lowest MAE, MSE and RMSE values among the six tested models. This indicates that tree-based gradient boosting can capture the non-linear structure of the BIST XKTUM dataset more effectively than the other models examined in this study. However, the portfolio optimisation results show that predictive accuracy alone is not sufficient for evaluating the practical usefulness of a model. While XGBoost performs best in stock price prediction, the portfolio optimisation stage focuses on LSTM, Bi-LSTM and GRU because these sequence-based models are designed to generate temporal forecasts that can be integrated into the mean-variance optimisation framework.

The distinction between prediction accuracy and portfolio performance is one of the central findings of this paper. The LSTM-based portfolio achieves the highest ending equity value, suggesting that it is effective in translating sequential forecasts into cumulative portfolio growth. However, the Bi-LSTM portfolio records the highest annualised Sharpe ratio, indicating the most favourable balance between return and risk. This result is important because Islamic portfolio management should not be assessed only through return maximisation. A Shariah-compliant investment strategy should also consider risk discipline, portfolio stability and ethical investment boundaries. Therefore, the Bi-LSTM result is particularly relevant for investors who prioritise risk-adjusted performance rather than maximum terminal wealth.

The findings also contribute to the discussion on Islamic fintech and data-driven investment management. Islamic finance is often criticised for reproducing conventional financial practices without sufficiently realising its broader ethical and social objectives. The results of this study suggest that computational methods can be used to strengthen the operational capacity of Islamic finance when they are applied within a Shariah-screened investment universe. In this study, the BIST XKTUM index provides the ethical investment boundary, while machine learning models provide predictive and optimisation support. This combination allows portfolio construction to remain aligned with Shariah-compliant screening while benefiting from modern analytical tools.

From a methodological perspective, the results show the value of comparing both sequence-based and non-sequence-based models. LSTM, Bi-LSTM and GRU are designed to capture temporal dependencies in financial time series, whereas XGBoost, Random Forest and SVM operate through different modelling assumptions. The superior prediction performance of XGBoost suggests that non-sequential models may perform strongly when the dataset contains engineered features that capture return, volatility and other market patterns. At the same time, the portfolio results show that recurrent models remain important when the objective is not only prediction, but also the generation of sequential return forecasts for portfolio allocation.

The portfolio return distributions provide further evidence that each model produces a different risk profile. The LSTM portfolio generates the highest ending equity, but with higher annualised volatility than Bi-LSTM and GRU. The Bi-LSTM portfolio produces a more balanced outcome, combining relatively strong ending equity with the highest Sharpe ratio and the lowest volatility. The GRU portfolio produces lower terminal equity, but remains competitive in terms of volatility and risk-adjusted performance. These results suggest that the choice of model should depend on the investor's objective. A growth-oriented investor may prefer LSTM, while a risk-conscious investor may prefer Bi-LSTM.

The findings also have implications for the development of Islamic fintech in Türkiye. The BIST XKTUM provides a relevant market setting for testing data-driven Shariah-compliant investment strategies. By combining Shariah screening, machine learning prediction and portfolio optimisation, this study demonstrates how Islamic fintech can move beyond payment systems and digital banking toward more advanced investment management applications. Such applications may support the development of automated Shariah-compliant advisory systems, risk-aware portfolio construction tools and data-driven investment platforms.

Overall, the results indicate that machine learning can support Shariah-compliant portfolio management, but model evaluation must be multidimensional. XGBoost is the strongest model for prediction accuracy, LSTM is the strongest model for ending equity and Bi-LSTM is the strongest model for risk-adjusted performance. This confirms that no single model dominates across all criteria. Therefore, future Islamic fintech applications should consider ensemble decision frameworks or multi-criteria model selection approaches that jointly evaluate forecasting error, cumulative portfolio growth, volatility and Sharpe ratio.